\subsection*{Locating a glyph when its grammatical role changes}
The earlier reader failed on all 14 object-first passages in the two-book classroom test. We annotated both ends of the verb--object relation on the training images, fitted a joint source-position reader, and trained a separate classifier to read which endpoint the pictured question requested. On 18 held pairs from four other books, the verb and its object were identical across the two passages but occurred in opposite orders. Given the intended role, both uniform training and oversampling of fronted objects located the two words on both pages in eight of 18 pairs. A model with exchanged endpoint labels reached nine. This test gave no evidence of an advantage for oversampling. Moreover, the new question classifier initially read both English classroom questions as requests for the object. All three neural fits consequently scored zero complete bilingual pairs.

We kept the source reader and writer fixed and added 54 reviewed question pairs from a local language model during training. The resulting visual model, running independently, completed four of the 18 opposite-order pairs, compared with zero for the earlier image reader and five for a training-count baseline given the intended role (Fig.~\ref{fig:visual-clause-roles}). On the 14 earlier object-first excerpts, five now received all four correct written answers. However, both the original classifier and its distilled successor chose both roles correctly on the same nine of twelve separately reserved question phrasings. Repair of the original classroom question did not improve this transfer count.

A second test matched 429 passage pairs in which the same printed glyph served as a verb in one passage and as an object with a different English gloss in the other. Under the published English and Chinese questions, the revised model located both occurrences in 226 pairs and located and wrote both registered meanings in 110. The earlier reader scored 112 and 49, and the glyph-count baseline scored 80 and 38. In two unmarked passages the output for the pictured character meaning speech switched from \textit{speak} to \textit{speech}. A new wording of the action question directed the reader to the wrong character and the reply became \textit{it}. None of the 429 pairs passed all twelve reserved phrasings with the distilled model.
